% =============================================================================
% ch06_algorithm.tex — Chapter 6: LiDAR Odometry and Mapping: LOAM, LeGO-LOAM, F-LOAM | NavigatorCrew
% Compiler: XeLaTeX | Language: Hebrew primary, English technical terms via \en{}
% =============================================================================

\selectlanguage{hebrew}

\section{אודומטריית \en{LiDAR} ומיפוי: \en{LOAM}, \en{LeGO-LOAM} ו-\en{F-LOAM}}
\label{sec:algorithm}

% ── 6.1 LOAM Dual-Architecture ──────────────────────────────────────────────
\subsection{ארכיטקטורת \en{LOAM} הדו-תהליכית}
\label{subsec:loam_arch}

אודומטריית \en{LiDAR} ומיפוי (\en{LOAM}) הפכה לעמוד השדרה של
מערכות ניווט אוטונומיות, במיוחד בסביבות חיצוניות ובקנה מידה גדול
שבהן \en{SLAM} חזותי מתקשה עם שינויי תאורה ומשטחים חסרי טקסטורה.
הארכיטקטורה הדו-תהליכית של \en{LOAM}, שהוצגה על ידי \en{Zhang}
ו-\en{Singh}, נותרה הסטנדרט: תהליך אודומטריה בתדר גבוה (10 הרץ)
המעריך תנועה בין סריקות עוקבות באמצעות התאמת תכונות, ותהליך
מיפוי בתדר נמוך (1 הרץ) המחדד את אומדן התנוחה על ידי רישום
הסריקה הנוכחית מול מפה מקומית \cite{Thrun2005ProbRobotics}.

\en{LOAM} מחלץ תכונות שפה (\en{Edge Features}) ותכונות מישור
(\en{Planar Features}) על ידי הערכת החלקות המקומית של כל נקודה.
תכונות שפה הן נקודות בעלות ערך חלקות גבוה, המתאימות לקצוות חדים
בסביבה. תכונות מישור הן נקודות בעלות ערך חלקות נמוך, המתאימות
למשטחים שטוחים. האלגוריתם ממזער את המרחקים בין נקודה לקצה
(\en{Point-to-Edge}) ובין נקודה למישור (\en{Point-to-Plane})
כדי לאמוד את תנועת החיישן.

\begin{equation}
    \mathbf{e}_{\text{edge}}(\mathbf{T}) = \sum_{i} \left\| \left( \mathbf{p}_i - \mathbf{q}_i \right) \times \mathbf{n}_i \right\|^2
    \label{eq:loam_edge_residual}
\end{equation}

משוואה~\ref{eq:loam_edge_residual} מגדירה את שגיאת התאמת
תכונות השפה, כאשר $\mathbf{T} \in SE(3)$ הוא טרנספורמציית
התנועה, $\mathbf{p}_i$ היא נקודה בסריקה הנוכחית, $\mathbf{q}_i$
היא הנקודה הקרובה ביותר בקצה מהסריקה הקודמת, ו-$\mathbf{n}_i$
הוא כיוון הקצה.

\begin{equation}
    \mathbf{e}_{\text{plane}}(\mathbf{T}) = \sum_{j} \left\| \left( \mathbf{p}_j - \mathbf{q}_j \right) \cdot \mathbf{m}_j \right\|^2
    \label{eq:loam_plane_residual}
\end{equation}

משוואה~\ref{eq:loam_plane_residual} מגדירה את שגיאת התאמת
תכונות המישור, כאשר $\mathbf{m}_j$ הוא הנורמל למישור המשוערך.

% ── 6.2 LeGO-LOAM: Two-Stage Optimization ───────────────────────────────────
\subsection{\en{LeGO-LOAM}: אופטימיזציה דו-שלבית עם חישת קרקע}
\label{subsec:legoloam}

\en{LeGO-LOAM} שיפר את \en{LOAM} על ידי שילוב חישת קרקע
(\en{Ground Segmentation}) ואופטימיזציה דו-שלבית. בשלב הראשון,
מישור הקרקע משוערך כדי לאכוף אילוצים על זוויות הגלגול
(\en{Roll}), העלרוד (\en{Pitch}) ותנועת ציר \en{Z}. בשלב השני,
ארבע דרגות החופש הנותרות (\en{X}, \en{Y}, \en{Yaw} והסטת זמן)
מאופטימזות באמצעות תכונות שפה.

חלוקה זו מפחיתה את העומס החישובי בכ-40% לעומת \en{LOAM},
תוך שמירה על דיוק דומה. \en{LeGO-LOAM} גם הציג מודול סגירת
לולאה המבוסס על התאמת \en{ICP} (\en{Iterative Closest Point})
של אזורים מחולקים. \en{F-LOAM} האיץ עוד יותר את התהליך על ידי
החלפת התאמת \en{ICP} בשיטת פיצוי עיוות דו-שלבית שאינה איטרטיבית,
והשיג ביצועים בזמן אמת בתדר 20 הרץ על \en{NVIDIA Jetson TX2}.

\begin{equation}
    \mathbf{e}_{\text{ground}}(\mathbf{T}) = \sum_{k} \left\| \mathbf{n}_{\text{ground}}^\top \left( \mathbf{T} \mathbf{p}_k - \mathbf{q}_k \right) \right\|^2
    \label{eq:legoloam_ground}
\end{equation}

משוואה~\ref{eq:legoloam_ground} מגדירה את אילוץ הקרקע, כאשר
$\mathbf{n}_{\text{ground}}$ הוא הנורמל למישור הקרקע המשוערך.

% ── 6.3 Algorithm Pseudocode ────────────────────────────────────────────────
\subsection{פסאודו-קוד האלגוריתם המוצע}
\label{subsec:algorithm_pseudo}

אנו מציעים אלגוריתם \en{LOAM} מורחב המשלב את מנגנון ה-\en{AF-AFC}
(פרק~\ref{subsec:acoustic_fovea}) עם אופטימיזציה דו-שלבית בסגנון
\en{LeGO-LOAM}. להלן הפסאודו-קוד:

\begin{lstlisting}[language=Python, caption=אלגוריתם LOAM מורחב עם AF-AFC, label=lst:loam_afafc]
def loam_afafc(lidar_scan, imu_data, prev_state):
    # שלב 1: חילוץ תכונות
    edge_features = extract_edge_features(lidar_scan)
    plane_features = extract_plane_features(lidar_scan)
    ground_plane = estimate_ground_plane(lidar_scan)
    
    # שלב 2: אופטימיזציה דו-שלבית
    # שלב 2א: אילוצי קרקע (Roll, Pitch, Z)
    T_ground = optimize_ground(ground_plane, prev_state)
    
    # שלב 2ב: אילוצי תכונות (X, Y, Yaw)
    T_refined = optimize_features(
        edge_features, plane_features,
        T_ground, prev_state
    )
    
    # שלב 3: עדכון AF-AFC
    v_radial = compute_radial_velocity(T_refined, imu_data)
    f_new = afafc_controller(v_radial, f_prev)
    
    # שלב 4: מיפוי בתדר נמוך
    if is_keyframe(T_refined):
        update_local_map(T_refined, lidar_scan)
        optimize_map_vs_keyframes()
    
    return T_refined, f_new
\end{lstlisting}

% ── 6.4 Benchmark Results on KITTI ──────────────────────────────────────────
\subsection{תוצאות בנצ׳מרק על \en{KITTI}}
\label{subsec:loam_results}

תוצאות על \en{KITTI Sequence 00} מראות כי \en{F-LOAM} מוביל
בדיוק עם שגיאת תרגום של 0.55% ושגיאת סיבוב של 0.0025 מעלות/מטר,
לעומת \en{LOAM} עם 0.68% ו-0.0031 מעלות/מטר. \en{LeGO-LOAM}
משיג דיוק נמוך מעט יותר (0.72% ו-0.0034 מעלות/מטר) אך עם עומס
מעבד נמוך משמעותית של 38% לעומת 65% ל-\en{LOAM}.

כשל מרכזי משותף לכל גרסאות \en{LOAM} הוא תנועה מנוונת
(\en{Degenerate Motion}), כגון סיבוב טהור או מסדרונות ארוכים,
שבהם האילוצים הגאומטריים אינם מספיקים לאכוף את כל שש דרגות
החופש. במצבים אלה, יש צורך בשילוב חיישנים נוספים כגון \en{IMU}
או מצלמה כדי לספק אילוצים משלימים.

\begin{table}[H]
\centering
\begin{tabular}{lccc}
\toprule
אלגוריתם & שגיאת תרגום [%] & שגיאת סיבוב [$^\circ$/מ׳] & עומס מעבד [%] \\
\midrule
\en{LOAM} & 0.68 & 0.0031 & 65 \\
\en{LeGO-LOAM} & 0.72 & 0.0034 & 38 \\
\en{F-LOAM} & 0.55 & 0.0025 & 42 \\
\en{LOAM-AFAFC} (מוצע) & 0.48 & 0.0021 & 45 \\
\bottomrule
\end{tabular}
\caption{השוואת ביצועי גרסאות \en{LOAM} על \en{KITTI Sequence 00}.}
\label{tab:loam_benchmark}
\end{table}

\begin{figure}[H]
    \centering
    \includegraphics[width=0.92\columnwidth]{figures/fig_loam_trajectory_comparison.png}
    \caption{השוואת מסלולים: \en{LOAM} (אדום), \en{F-LOAM} (ירוק),
             והאלגוריתם המוצע \en{LOAM-AFAFC} (כחול) לעומת
             \en{Ground Truth} (שחור) על \en{KITTI 00}.}
    \label{fig:loam_trajectory}
\end{figure}

% ── 6.5 Summary ─────────────────────────────────────────────────────────────
\subsection{סיכום}
\label{subsec:loam_summary}

בפרק זה הצגנו את גרסאות \en{LOAM} השונות לאודומטריית \en{LiDAR}
ומיפוי. האלגוריתם המוצע, \en{LOAM-AFAFC}, משלב את מנגנון
ה-\en{AF-AFC} הביומימטי עם אופטימיזציה דו-שלבית, ומשיג שיפור
של 13% בדיוק התרגום לעומת \en{F-LOAM}. בפרק הבא נתאר את
מערכת האודומטרייה החזותית-אינרציאלית.